%% ============================================================
%% CV — Ma, Ruiyang
%% Source of truth: ../CONTENT.md
%% Compile: pdflatex cv.tex (run twice)
%% ============================================================
\documentclass[11pt, letterpaper]{article}

%% ---- Packages -----------------------------------------------
\usepackage[T1]{fontenc}
\usepackage[utf8]{inputenc}
\usepackage{lmodern}
\usepackage[
  top    = 0.75in,
  bottom = 0.75in,
  left   = 0.85in,
  right  = 0.85in
]{geometry}
\usepackage{xcolor}
\usepackage{hyperref}
\usepackage{microtype}
\usepackage{enumitem}
\usepackage{array}
\usepackage{booktabs}
\usepackage{parskip}
\usepackage{fontawesome5}
\usepackage{tabularx}
\usepackage{calc}
\usepackage{etoolbox}

%% ---- Colors -------------------------------------------------
\definecolor{accent}{RGB}{30, 64, 175}    % deep blue (matches website)
\definecolor{muted}{RGB}{100, 116, 139}   % slate-500

%% ---- Hyperref -----------------------------------------------
\hypersetup{
  colorlinks = true,
  urlcolor   = accent,
  linkcolor  = accent,
}

%% ---- Typography tweaks --------------------------------------
\setlength{\parindent}{0pt}
\setlength{\parskip}{0pt}
\setlength{\tabcolsep}{0pt}

%% ---- Section heading ----------------------------------------
% Prints: SECTION TITLE with a thin colored rule below
\newcommand{\cvsection}[1]{%
  \vspace{14pt}%
  {\color{accent}\rule{\linewidth}{1.8pt}}\\[-6pt]
  {\large\bfseries\color{accent} #1}%
  \vspace{5pt}%
}

%% ---- Dated entry: \entry{DATE}{BODY} -----------------------
% Date in left col (1.55in), body in right col
\newlength{\datecol}
\setlength{\datecol}{1.55in}
\newcommand{\entry}[2]{%
  \begin{tabular}{@{}p{\datecol}@{\hspace{6pt}}p{\linewidth-\datecol-6pt}@{}}
    \raggedleft\small\color{muted} #1 &
    \raggedright #2 \\
  \end{tabular}%
  \vspace{3pt}%
}

%% ---- Paper entry -------------------------------------------
\newenvironment{paperentry}{%
  \vspace{2pt}%
  \begin{list}{}{%
    \setlength{\leftmargin}{0pt}%
    \setlength{\topsep}{0pt}%
    \setlength{\itemsep}{0pt}%
    \setlength{\parsep}{3pt}%
  }\item%
}{%
  \end{list}\vspace{5pt}%
}

%% ============================================================
\begin{document}
\thispagestyle{empty}
\pagestyle{plain}

%% ---- HEADER ------------------------------------------------
\begin{center}
  {\fontsize{26}{30}\selectfont\bfseries Ma, Ruiyang}\\[8pt]
  \color{muted}
  \faEnvelope[regular]\enspace
    \href{mailto:rma13@illinois.edu}{\textcolor{accent}{rma13@illinois.edu}}
  \quad\textbullet\quad
  \faGlobe\enspace
    \href{https://ruiyangma.github.io/ruiyang-ma.github.io/}%
         {\textcolor{accent}{ruiyangma.github.io}}
  \quad\textbullet\quad
  \faUniversity\enspace
    University of Illinois at Urbana-Champaign
\end{center}

\vspace{2pt}
{\color{accent}\rule{\linewidth}{0.4pt}}
\vspace{4pt}

%% ---- RESEARCH PROFILE --------------------------------------
\cvsection{Research Profile}

\begin{tabular}{@{}p{\datecol}@{\hspace{6pt}}p{\linewidth-\datecol-6pt}@{}}
  \small\itshape\color{muted}Fields &
    Macro Finance;\enspace Monetary Economics \\[4pt]
  \small\itshape\color{muted}Topics &
    Monetary Transmission \textbullet\enspace
    Networks \textbullet\enspace
    Financial Frictions \textbullet\enspace
    Identification in Macroeconomics \\
\end{tabular}

%% ---- EDUCATION ---------------------------------------------
\cvsection{Education}

\entry{2023 -- Present}{%
  \textbf{Ph.D.} in Economics\\
  University of Illinois at Urbana-Champaign}

\entry{2021 -- 2023}{%
  \textbf{M.Sc.} in Policy Economy\\
  University of Illinois at Urbana-Champaign}

\entry{2018 -- 2020}{%
  \textbf{B.A.} in Economics\\
  University of Illinois at Urbana-Champaign}

\entry{2016 -- 2018}{%
  \textbf{B.S.} in Economics and Mathematics\\
  Southwestern University of Finance and Economics}

%% ---- RESEARCH ----------------------------------------------
\cvsection{Working Papers}

%% Paper 01
\begin{paperentry}
  \textbf{The Financial Reallocation Channel of Monetary Policy}
  \hfill{\small\itshape\color{muted}Work in Progress}\\[2pt]
  {\small\color{muted}with Junze Lin}\\[5pt]
  We embed a bank dependency network---a matrix whose entries measure the marginal
  impact of one bank's fundamentals on another's market value---into a New Keynesian
  model with heterogeneous households. The network generates heterogeneous bank
  responses to a common policy impulse, creating a financial reallocation channel that
  simultaneously stimulates aggregate demand and expands aggregate supply following an
  expansionary monetary shock, reinforcing each other on output but partially
  offsetting each other on inflation, thereby endogenously flattening the Phillips
  Curve. The equilibrium is governed by two network objects---a credit multiplier
  matrix that amplifies contemporaneous shocks and a legacy matrix that propagates the
  credit distribution across time---making the effective Phillips Curve slope
  state-dependent on the history of prior policy. Quantitative easing is amplified
  strictly more than conventional interest-rate policy because it transmits entirely
  through bank balance sheets.
\end{paperentry}

%% Paper 03
\begin{paperentry}
  \textbf{The Financial Reallocation Channel: How Interbank Networks Shape Spatial
  Misallocation}
  \hfill{\small\itshape\color{muted}Work in Progress}\\[2pt]
  {\small\color{muted}with Ke Gao and Hongrui Pu}\\[5pt]
  We propose that the topology of interbank networks shapes spatial allocative
  efficiency through a financial reallocation channel. Banks with high
  self-amplification centrality face greater systemic risk exposure, leading them to
  contract local working-capital lending. This tighter credit supply generates
  endogenous microeconomic input wedges for local firms. Using the 1863--64 National
  Banking Acts as an institutional shock that rewired the network, we demonstrate that
  regions whose banks move to less exposed positions experience relaxed credit
  constraints. Because this spatial expansion disproportionately benefits financially
  constrained, inefficiently small firms, it corrects pre-existing misallocation and
  boosts aggregate efficiency. Furthermore, the layered structure of interbank claims
  amplifies these local credit distortions. Ultimately, we establish that financial
  network architecture and trade integration jointly determine aggregate macroeconomic
  productivity.
\end{paperentry}

%% Paper 04
\begin{paperentry}
  \textbf{Pricing the Risk Lock-in: Financial Frictions, Endogenous Production
  Networks, and Corporate Bond Spreads}
  \hfill{\small\itshape\color{muted}Work in Progress}\\[2pt]
  {\small\color{muted}with Shaoqing Pan and Rong Yuwen}\\[5pt]
  This research proposes to investigate how financial frictions distort firm-level
  supply chain adjustments and the resulting credit risk pricing. We develop a
  theoretical framework to identify a ``risk lock-in'' mechanism, where firms facing
  financing constraints and fixed adjustment costs may be unable to reconfigure their
  supplier networks toward stability during periods of high uncertainty. The study
  examines how these constraints break the ``flight-to-stability'' motive, potentially
  creating a ``volatility trap'' that depresses firm productivity, and how the market
  prices this lack of flexibility, leading to a divergence in credit spreads between
  constrained and unconstrained firms.
\end{paperentry}

%% ---- TEACHING ----------------------------------------------
\cvsection{Teaching Experience}

{\small\itshape\color{muted}University of Illinois at Urbana-Champaign}

\vspace{6pt}

\entry{FA25}{%
  Intermediate Macroeconomics \enspace\textbullet\enspace
  Undergraduate \hfill \textbf{Head Teaching Assistant}}

\entry{SP26, SP25, FA24}{%
  Microeconomics Principles \enspace\textbullet\enspace
  Undergraduate \hfill Teaching Assistant}

\entry{SP24}{%
  Macroeconomics II \enspace\textbullet\enspace
  Ph.D. \hfill Teaching Assistant}

\entry{FA23}{%
  Intermediate Macroeconomics \enspace\textbullet\enspace
  Undergraduate \hfill Teaching Assistant}

%% ---- AWARDS ------------------------------------------------
\cvsection{Awards \& Fellowships}

\entry{2026}{Paul W. Boltz Fellowship, University of Illinois}

\entry{2025}{Hung-Chao and Julia Tai Family Fellowship, University of Illinois}

%% ---- SKILLS ------------------------------------------------
\cvsection{Skills}

\begin{tabular}{@{}p{\datecol}@{\hspace{6pt}}p{\linewidth-\datecol-6pt}@{}}
  \small\itshape\color{muted}Languages &
    Chinese (native)\enspace\textbullet\enspace English (fluent) \\[4pt]
  \small\itshape\color{muted}Computing &
    Python \enspace\textbullet\enspace
    Dynare \enspace\textbullet\enspace
    Stata \enspace\textbullet\enspace
    \LaTeX \\
\end{tabular}

\vspace{12pt}
{\color{muted}\rule{\linewidth}{0.4pt}}\\[3pt]
{\small\color{muted}\itshape Updated: \today}

\end{document}
